% Chapter 4: MHV Amplitudes
% Based on arXiv:2602.12176 and research from 2026-02-17

\chapter{Maximally Helicity Violating (MHV) Amplitudes}
\label{ch:mhv}

\section{Why ``Maximally Helicity Violating''?}

For a tree-level $n$-gluon amplitude $A_n(h_1,\dots,h_n)$ with $h_i=\pm 1$:
\begin{itemize}
  \item Let $m_-$ denote the number of negative-helicity gluons.
  \item Naive Yang--Mills vertices with all same helicity ($+\cdots+$ or $-\cdots-$) are strongly suppressed --- amplitudes with $m_-=0$ or $m_-=n$ vanish at tree level (and at all loops).
  \item Amplitudes with $m_-=1$ (single-minus) were long conjectured to vanish; see \S\ref{sec:single-minus}.
  \item In twistor space (Witten 2004), tree amplitudes are supported on curves of degree $d = m_- - 1 + \ell$, so at tree level $\ell=0$ the non-zero sector requires $0 \le m_- \le 4$ generically.
\end{itemize}
\textbf{MHV = maximally helicity violating} corresponds to $m_- = 2$, i.e.\ $m_+ = n-2$. This is the maximum number of ``wrong-helicity'' gluons that gives a generically non-zero amplitude, in the sense that every cyclic ordering is non-zero. The name reflects that having only 2 negative-helicity gluons out of $n$ is the extreme violation of a naive helicity-conservation intuition.

\section{The Parke-Taylor Formula}

For $n$-gluon scattering with legs $j$ and $k$ carrying negative helicity, the tree-level color-ordered amplitude is:
\begin{equation}
\boxed{A_n^{\text{MHV}}(\ldots, j^-, \ldots, k^-, \ldots) = i \frac{\langle j k \rangle^4}{\langle 12 \rangle \langle 23 \rangle \cdots \langle n1 \rangle}}
\end{equation}
(momentum-conserving delta function $\delta^4(P)$ implicit).
For the canonical choice $(j,k)=(1,2)$ this reads $i\,\langle 12\rangle^4 / (\langle 12\rangle\langle 23\rangle\cdots\langle n1\rangle)$.

This compact formula replaces $O(n!)$ Feynman diagrams with a single expression.

\paragraph{Little group weight check.}
Under $\lambda_i \to t_i \lambda_i$, $\tilde\lambda_i \to t_i^{-1}\tilde\lambda_i$, a helicity-$h_i$ gluon requires the amplitude to scale as $t_i^{-2h_i}$.
\begin{itemize}
  \item Negative-helicity legs $j,k$: numerator $\langle jk\rangle^4 \sim t_j^4 t_k^4$, denom contributes $t_j^{-1}t_k^{-1}$ (one factor each), net $t_j^4/t_j^{-1}=t_j^{4+1}$... more carefully: $\langle jk\rangle^4$ in numerator gives weight $+4$ for $j$ and $+4$ for $k$; each $\langle \cdot j\cdot\rangle$ in denominator gives $-1$; there are exactly 2 such factors, so net weight for $j$ is $4-2=+2 = -2h_j$ with $h_j=-1$. \checkmark
  \item Positive-helicity legs: appear twice in the denominator (once as $\langle i{,}i{+}1\rangle$ and once as $\langle i{-}1{,}i\rangle$), giving weight $-2 = -2(+1)$. \checkmark
\end{itemize}

\section{Helicity Configuration Table}

\begin{center}
\begin{tabular}{|c|c|c|}
\hline
Helicity Configuration & Tree Amplitude & Status \\
\hline
$(-,-,+,+,\ldots,+)$ & $\neq 0$ & MHV \\
$(-,+,+,+,\ldots,+)$ & $\neq 0$ & MHV (Parke-Taylor) \\
$(+,-,+,+,\ldots,+)$ & $\neq 0$ & MHV (cyclic) \\
$(-,-,-,+,\ldots,+)$ & $\neq 0$ & NMHV \\
$(-,+,+,+,+,+)$ & $\neq 0$ & Single-minus (2026 result) \\
$(+,+,+,+,+,+)$ & $= 0$ & Vanishing \\
\hline
\end{tabular}
\end{center}

\section{Single-Minus Amplitudes Are Nonzero}
\label{sec:single-minus}

Recent work (Guevara, Lupsasca, Skinner, Strominger, Weil, arXiv:2602.12176) proves that single-minus ($m_-=1$) tree amplitudes are non-zero for all $n\ge 4$.

\begin{equation}
A_n(1^-, 2^+, \ldots, n^+) \neq 0 \quad \forall\, n \ge 4.
\end{equation}

This breaks the longstanding conjecture that only $m_-\ge 2$ configurations were generically nonvanishing.

\paragraph{Half-collinear locus (arXiv:2602.12176 Eq.~(13)).}
The single-minus amplitude is supported on the \emph{half-collinear locus}:
\begin{equation}
\langle ij\rangle = 0 \quad \forall\, i,j \in \{1,\dots,n\}.
\label{eq:half-collinear-locus}
\end{equation}
This means all $\lambda_i$ are proportional, i.e.\ $\lambda_i = c_i\,\lambda_1$ for some $c_i\in\mathbb{C}$, so all $n$ momenta are collinear. The $\delta(\langle 1a\rangle)$ factors in the ansatz below enforce this support explicitly for $a=2,\dots,n$.

\paragraph{Explicit single-minus ansatz (arXiv:2602.12176 Eq.~(15)).}
The ansatz is:
\begin{align}
\mathcal{A}_n(1^-,2^+,\dots,n^+)
&= i^{2-n}\, \frac{\langle r 1\rangle^{n+1}}{\langle r 2\rangle \langle r 3\rangle \cdots \langle rn\rangle}\, A_{1\cdots n}
\prod_{a=2}^{n} \delta\!\bigl(\langle 1 a\rangle\bigr) \notag\\
&\quad \times\; \delta^2\!\Bigl(\sum_{i=1}^n \langle r i\rangle\, \tilde\lambda_i\Bigr)
\label{eq:single-minus-ansatz}
\end{align}
where $r$ is the reference spinor, $A_{1\cdots n}$ is a scalar prefactor, and $\delta^2(\cdots)$ is a 2-component ($\dot\alpha$) delta function imposing a residual momentum constraint. In affine coordinates (\S\ref{sec:affine-spinors}), $\langle 1a\rangle = z_a - z_1$, so the delta functions set $z_a = z_1$ for all $a$.

\paragraph{MHV comparison (arXiv:2602.12176 Eq.~(11)).}
For comparison, the MHV amplitude in the same notation with negative-helicity legs $r,s$ reads:
\begin{equation}
\mathcal{A}_n^{\text{MHV}} = i\,\langle rs\rangle^4\;\mathrm{PT}_{\mathrm{cyc}}\;\delta^4\!\Bigl(\sum_{k=1}^n p_k\Bigr),
\label{eq:mhv-ptcyc}
\end{equation}
where $\mathrm{PT}_{\mathrm{cyc}} \equiv 1/(\langle 12\rangle\langle 23\rangle\cdots\langle n1\rangle)$ is the cyclic Parke-Taylor denominator. The full 4-momentum-conserving delta function $\delta^4(\sum_k p_k)$ is displayed explicitly here (it was left implicit in \eqref{eq:mhv-ptcyc}). The structural contrast with \eqref{eq:single-minus-ansatz} is striking: MHV has a standard $\delta^4(P)$ on generic kinematics, while single-minus is localized on the half-collinear locus \eqref{eq:half-collinear-locus} via the $\prod_a\delta(\langle 1a\rangle)$ factors.

\section{BCJ Relations}

The Bern-Carrasco-Johansson (BCJ) relations express amplitude redundancies:
\begin{equation}
\sum_{i=2}^{n-1} \frac{s_{1i}}{\langle 1i \rangle [i1]} A_n(2, \ldots, i, 1, i+1, \ldots, n) = 0
\end{equation}

These reduce the number of independent color-ordered amplitudes from $(n-3)!$ to $(n-3)!/2$.

\section{Double Copy}

Gauge theory amplitudes factorize into gravity amplitudes via:
\begin{equation}
M_n^{\text{gravity}} = \sum_{\text{duals}} \frac{A_n^{\text{gauge}} \tilde{A}_n^{\text{gauge}}}{\text{kinematic factors}}
\end{equation}

This is the KLT (Kawai-Lewellen-Tye) relation.
